\chapter{Equilibration Estimation and Ablation-Threshold Search}
\label{app:equilibration-ablation}

\section{Equilibration index estimation}
\label{app:equilibration}

Let the surface equilibration signal be defined by the electron--lattice
temperature gap
\begin{equation}
    s(t) \coloneqq T_e(\mathbf{r}_{\mathrm{surf}}, t) - T_l(\mathbf{r}_{\mathrm{surf}}, t),
    \label{eq:app-signal}
\end{equation}
where \(\mathbf{r}_{\mathrm{surf}}\) is a prescribed representative point on the irradiated surface.
In the one-dimensional setting, \(\mathbf{r}_{\mathrm{surf}}\) is the surface node \(z=z_s\).
In the two-dimensional setting, one may select any maximiser of the peak surface electron temperature,
\begin{equation}
    \varrho_* \in \operatorname*{arg\,max}_{\varrho}
    \max_{0 \le t \le \tau_f} T_e(\varrho,z_s,t),
    \qquad
    \mathbf{r}_{\mathrm{surf}} \coloneqq (\varrho_*,z_s),
    \label{eq:app-surface-argmax}
\end{equation}
with the understanding that if the maximiser is not unique, any fixed choice is admissible.

Let \(t_n\), \(n=0,\dots,N-1\), denote the discrete output times and define
\begin{equation}
    s_n \coloneqq s(t_n).
\end{equation}
If optional smoothing is applied, let \(w\ge 1\) be an odd window width and define the truncated centred moving average
\begin{equation}
    I_n \coloneqq
    \left\{
    m\in\mathbb{Z} :
    |m|\le \frac{w-1}{2},
    \ 0\le n-m \le N-1
    \right\},
    \qquad
    \tilde s_n \coloneqq \frac{1}{|I_n|}\sum_{m\in I_n} s_{n-m}.
    \label{eq:app-smoothed-series}
\end{equation}
When no smoothing is required, one may simply set \(\tilde s_n \equiv s_n\).
This role of local smoothing is consistent with the Savitzky--Golay viewpoint of local least-squares filtering, although the present implementation uses only a moving average unless higher-order smoothing is explicitly desired \cite{Savitzky_1964}.

Define the peak index by the first maximiser of the processed sequence,
\begin{equation}
    n_{\mathrm{peak}}
    \coloneqq
    \min\operatorname*{arg\,max}_{0\le n < N} \tilde s_n,
    \qquad
    s_{\mathrm{peak}} \coloneqq \tilde s_{n_{\mathrm{peak}}}.
    \label{eq:app-peak-index}
\end{equation}
The pre-pulse baseline may be taken as
\begin{equation}
    s_{\mathrm{base}} \coloneqq \tilde s_0,
    \qquad
    \Delta_s \coloneqq \max\{s_{\mathrm{peak}}-s_{\mathrm{base}},\,0\}.
    \label{eq:app-baseline-rise}
\end{equation}
If several pre-pulse output samples are available, \(s_{\mathrm{base}}\) may alternatively be replaced by their average.

Introduce an absolute tolerance \(\varepsilon_{\mathrm{abs}}>0\) (K) and a relative tolerance
\(\varepsilon_{\mathrm{rel}}\in(0,1)\). The adaptive return threshold is
\begin{equation}
    \Theta
    \coloneqq
    s_{\mathrm{base}}
    +
    \max\!\bigl(
    \varepsilon_{\mathrm{abs}},
    \varepsilon_{\mathrm{rel}}\Delta_s
    \bigr).
    \label{eq:app-threshold}
\end{equation}
This construction is analogous in spirit to adaptive-threshold rules based on signal and noise levels, as in the Pan--Tompkins framework, though here the threshold is referenced to the baseline and peak excursion of the equilibration signal rather than to running electrocardiogram peak statistics \cite{Pan_1985}.

The equilibration index is then defined by the first post-peak threshold return,
\begin{equation}
    n_{\mathrm{eq}}
    \coloneqq
    \begin{cases}
        n_{\mathrm{peak}},
         & \Delta_s \le \varepsilon_{\mathrm{abs}}, \\[0.5ex]
        \min\{n\ge n_{\mathrm{peak}} : \tilde s_n \le \Theta\},
         & \text{if such an index exists},          \\[0.5ex]
        N-1,
         & \text{otherwise}.
    \end{cases}
    \label{eq:app-neq}
\end{equation}
The associated physical equilibration-time estimator is
\begin{equation}
    \hat t_{\mathrm{eq}} \coloneqq t_{n_{\mathrm{eq}}}.
    \label{eq:app-teq}
\end{equation}

This is a single-threshold post-peak crossing rule.
It is therefore not a hysteresis detector in the strict sense of Canny,
where a high threshold starts a structure and a lower threshold continues it under a connectivity condition \cite{Canny_1986}.
Its role here is instead to suppress premature post-peak returns caused by mild oscillations or temporal discretisation effects.


\section{Exponential bracketing and bisection refinement}
\label{app:exp-bisection}

\subsection{Threshold definition and structural assumptions}

For fixed material parameters and fixed pulse parameters, define
\begin{equation}
    T_{l,\max}(F)
    \coloneqq
    \max_{(t,\mathbf{r})\in(0,\tau_f]\times\Omega}
    T_l(\mathbf{r},t;F),
    \label{eq:app-tlmax}
\end{equation}
and the melting-onset predicate
\begin{equation}
    D(F)
    \iff
    T_{l,\max}(F)\ge T_m.
    \label{eq:app-damage-predicate}
\end{equation}
The threshold fluence is
\begin{equation}
    F_{\mathrm{th}}
    \coloneqq
    \inf\{F>0 : D(F)\}.
    \label{eq:app-threshold-fluence}
\end{equation}

We assume:
\begin{enumerate}
    \item \(D(0)=\mathrm{false}\);
    \item \(D\) is monotone nondecreasing, i.e.
          \[
              F_a<F_b \Longrightarrow D(F_a)\Rightarrow D(F_b);
          \]
    \item when a sharp complexity count is desired, \(T_{l,\max}(F)\) is continuous in \(F\), so that the threshold set is right-closed.
\end{enumerate}

\subsection{Exponential bracketing}

Let \(F_{\mathrm{init}}>0\) and \(r>1\). Define
\begin{equation}
    F_n \coloneqq F_{\mathrm{init}}r^n,
    \qquad n\in\mathbb{N}_0.
    \label{eq:app-exp-sequence}
\end{equation}
If the search succeeds, let
\begin{equation}
    n_* \coloneqq \min\{n\ge 0 : D(F_n)\}.
    \label{eq:app-exp-index}
\end{equation}
The initial bracket is
\begin{equation}
    F_l^{(0)} \coloneqq
    \begin{cases}
        0,         & n_*=0,    \\
        F_{n_*-1}, & n_*\ge 1,
    \end{cases}
    \qquad
    F_u^{(0)} \coloneqq F_{n_*}.
    \label{eq:app-initial-bracket}
\end{equation}
Under the monotonicity assumption, this implies
\[
    D(F_l^{(0)})=\mathrm{false},
    \qquad
    D(F_u^{(0)})=\mathrm{true},
    \qquad
    F_{\mathrm{th}}\in[F_l^{(0)},F_u^{(0)}].
\]

If \(F_{\mathrm{init}}\le F_{\mathrm{th}}\), then the first successful index satisfies the elementary bound
\begin{equation}
    n_*
    \le
    1+
    \left\lceil
    \frac{\log(F_{\mathrm{th}}/F_{\mathrm{init}})}{\log r}
    \right\rceil.
    \label{eq:app-exp-cost}
\end{equation}
If, in addition, \(D(F_{\mathrm{th}})=\mathrm{true}\) (for instance by continuity of \(T_{l,\max}\) in \(F\)), then the sharper identity holds:
\begin{equation}
    n_*
    =
    \left\lceil
    \frac{\log(F_{\mathrm{th}}/F_{\mathrm{init}})}{\log r}
    \right\rceil.
    \label{eq:app-exp-cost-sharp}
\end{equation}
The number of simulator evaluations in the exponential phase is therefore
\begin{equation}
    n_{\exp}=n_*+1.
\end{equation}

The relevance of \cite{Beigel_1990} is conceptual: it places this exponential-bracketing phase in the broader class of unbounded-search algorithms. In the present geometric schedule, however, \eqref{eq:app-exp-cost} is obtained directly from the inequality \(F_{\mathrm{init}}r^{n_*}\ge F_{\mathrm{th}}\).

\subsection{Bisection refinement}

Given a bracket \([F_l^{(0)},F_u^{(0)}]\) with
\[
    D(F_l^{(0)})=\mathrm{false},
    \qquad
    D(F_u^{(0)})=\mathrm{true},
\]
define
\begin{equation}
    F_m^{(k)}
    \coloneqq
    \frac{F_l^{(k)}+F_u^{(k)}}{2}.
    \label{eq:app-midpoint}
\end{equation}
The update is
\begin{equation}
    (F_l^{(k+1)},F_u^{(k+1)})
    \coloneqq
    \begin{cases}
        (F_l^{(k)},F_m^{(k)}), & D(F_m^{(k)})=\mathrm{true},  \\
        (F_m^{(k)},F_u^{(k)}), & D(F_m^{(k)})=\mathrm{false}.
    \end{cases}
    \label{eq:app-bisection-update}
\end{equation}
By monotonicity, the invariant
\[
    F_{\mathrm{th}}\in[F_l^{(k)},F_u^{(k)}]
\]
is preserved for all \(k\ge 0\).

The bracket width satisfies
\begin{equation}
    W_k \coloneqq F_u^{(k)}-F_l^{(k)}
    =
    \frac{F_u^{(0)}-F_l^{(0)}}{2^k},
    \label{eq:app-bisection-width}
\end{equation}
so it is sufficient to choose
\begin{equation}
    k
    \ge
    \left\lceil
    \log_2\!\left(
    \frac{F_u^{(0)}-F_l^{(0)}}{\varepsilon}
    \right)
    \right\rceil
    \label{eq:app-bisection-cost}
\end{equation}
in order to guarantee \(W_k\le \varepsilon\).

At termination, define the midpoint estimator
\begin{equation}
    \hat F_{\mathrm{th}}
    \coloneqq
    \frac{F_l+F_u}{2}.
    \label{eq:app-midpoint-estimate}
\end{equation}
Since \(F_{\mathrm{th}}\in[F_l,F_u]\), its deterministic enclosure error satisfies
\begin{equation}
    \bigl|\hat F_{\mathrm{th}}-F_{\mathrm{th}}\bigr|
    \le
    \delta F_{\mathrm{bis}}
    \coloneqq
    \frac{F_u-F_l}{2}.
    \label{eq:app-bisection-uncertainty}
\end{equation}

\subsection{Slope-propagated and combined uncertainty}

Suppose a local temperature uncertainty \(\delta T\) is available and that
\(T_{l,\max}\) is locally differentiable in \(F\) near the threshold.
Using two nearby fluence values \(F_a\neq F_b\), define the secant slope
\begin{equation}
    \hat S
    \coloneqq
    \frac{T_{l,\max}(F_a)-T_{l,\max}(F_b)}{F_a-F_b}.
    \label{eq:app-slope}
\end{equation}
If \(\hat S\neq 0\), the corresponding first-order propagated fluence uncertainty is
\begin{equation}
    \delta F_T
    \coloneqq
    \frac{\delta T}{|\hat S|}.
    \label{eq:app-temp-uncertainty}
\end{equation}

The bisection half-width \(\delta F_{\mathrm{bis}}\) is a deterministic enclosure.
If \(\delta F_T\) and \(\delta F_{\mathrm{other}}\) are instead interpreted as independent
standard uncertainties, then the combined standard uncertainty may be written as
\begin{equation}
    u_F
    \coloneqq
    \sqrt{
        \delta F_{\mathrm{bis}}^2
        +
        \delta F_T^2
        +
        \delta F_{\mathrm{other}}^2
    }.
    \label{eq:app-combined-uncertainty}
\end{equation}
If a conservative worst-case bound is preferred, one may use
\begin{equation}
    \delta F_{\mathrm{wc}}
    \coloneqq
    \delta F_{\mathrm{bis}}
    +
    \delta F_T
    +
    \delta F_{\mathrm{other}}.
\end{equation}
If no additional temperature-error estimate is available, the natural reported quantity is simply the deterministic bracketing uncertainty \(\delta F_{\mathrm{bis}}\).
